% =====================================================================
\section{Introduction}
% =====================================================================

\subsection{Context and motivation}

Strawberry production is among the most labour-intensive branches of
horticulture. Harvesting alone accounts for a large share of the
operating cost of a table-top greenhouse, it must be repeated every two
to three days through the season, and it is increasingly constrained by
the availability of seasonal labour. The task is also unusually hard to
automate: the fruit is soft, it grows in clusters where ripe and unripe
berries touch, it is partly hidden by leaves, and the picking window is
short. These properties have made strawberry harvesting a recurring
benchmark for agricultural robotics rather than a solved problem
\cite{bac2014harvesting,vanhenten2002cucumber}.

A greenhouse, on the other hand, is a far friendlier environment for a
mobile robot than an open field. It is flat, indoor, geometrically
regular, free of weather, and its layout --- parallel raised gutters
separated by service aisles --- is known in advance and does not change
between seasons. That combination is what makes a \emph{digital twin}
attractive: because the physical structure is stable and measurable, a
simulated greenhouse can be made dimensionally faithful, and a control
stack developed against it has a realistic chance of transferring to the
real installation.

\subsection{Objective of this work}

The objective of this Projet de Fin d'Ann\'ee is to build, from nothing,
a complete autonomous stack that lets a KUKA YouBot omnidirectional
mobile manipulator:

\begin{enumerate}[leftmargin=*,itemsep=1pt,topsep=2pt]
  \item build a metric map of an unknown greenhouse from a 2D lidar;
  \item localise itself in that map while it moves;
  \item plan and follow collision-free paths along the service aisles;
  \item detect ripe strawberries with an on-board colour camera and
        place them on the map;
  \item drive its five-DOF arm to pick them; and
  \item repeat the cycle autonomously, delivering the fruit to a depot.
\end{enumerate}

A secondary objective, imposed by the intended reuse of the platform in
healthcare logistics and indoor transport, is that the navigation,
mapping and safety layers must be \emph{domain-independent}: nothing in
them may assume strawberries, gutters, or this particular greenhouse.

\subsection{Two phases, and why}

The work is presented chronologically because it was carried out
chronologically, and because the second phase is only intelligible as a
consequence of the first.

\textbf{\phaseA{} (Webots).} A complete prototype was first written for
Webots R2025a, in plain Python, with no middleware. This was deliberate:
Webots supplies the robot model, the physics and the sensors out of the
box, so the entire effort could go into the algorithms themselves ---
mecanum kinematics, odometry, occupancy mapping, A$^\star$, pure
pursuit, colour vision, arm sequencing, and a behaviour-tree mission
supervisor. Fifteen controllers were written, ten of them single-purpose
test controllers used to validate one subsystem at a time. Phase~A
produced 3\,844 lines of Python and, more importantly, an algorithm
layer that had been exercised against a working simulator.

\textbf{\phaseB{} (ROS~2 / Gazebo).} The validated algorithms were then
lifted, essentially unmodified, into a framework-independent library
(\filename{youbot\_control/lib/}) and wrapped in ROS~2 Jazzy nodes, with
the scene rebuilt as a Gazebo Harmonic digital twin measured against the
real greenhouse. Phase~B is therefore not a rewrite but a
\emph{re-architecture}: the same A$^\star$, the same pure pursuit, the
same log-odds grid, now distributed across processes, with the problems
that distribution creates --- and with a measurement layer added so that
those problems become visible instead of being guessed at.

\subsection{What this report claims, and what it does not}

This report is written to be checked. Every quantitative statement in it
is traceable to a specific line of a specific simulation log, and the
procedure to regenerate those logs is given in Appendix~\ref{app:reproduce}.
Where a subsystem does not work, that is stated with the same emphasis
as where one does. In particular:

\begin{itemize}[leftmargin=*,itemsep=1pt,topsep=2pt]
  \item The mapping subsystem works well. The reconstructed interior is
        accurate to \measured{$-4$~cm~/~$+10$~cm} against a
        $9.90 \times 4.90$~m reference, with \measured{100\%} of the
        free floor observed and \measured{0.0\%} of it wrongly marked as
        obstacle (Sec.~\ref{sec:results}).
  \item Odometry calibration works. A one-off UMBmark-style calibration
        \cite{borenstein1996umbmark} holds position error at
        \measured{0.20~m} over a full mission, where the uncalibrated
        odometry diverges past \measured{10~m}.
  \item \emph{The homemade SLAM module does not work.} Incremental
        correlative scan matching without loop closure was implemented,
        instrumented, and measured to be \emph{worse} than the odometry
        prior it corrects. Sec.~\ref{sec:slam} explains why this is a
        property of the algorithm class and not a coding defect, and the
        production configuration disables it.
  \item \emph{The fruit detector is not accurate enough.} Colour
        thresholding reaches \measured{77\%} recall but only about one
        correct estimate in three, at \measured{0.22~m} localisation
        error --- roughly a quarter of the plant spacing, which means it
        identifies the right plant but not the right berry.
\end{itemize}

Reporting the negative results in full is not modesty. The instrumented
comparison between a homemade SLAM front-end and a calibrated dead-reckoning
baseline is one of the more useful outcomes of the project, and it only
exists because the measurement layer was built before the conclusions
were drawn.

\subsection{Contributions}

\begin{enumerate}[leftmargin=*,itemsep=1pt,topsep=2pt]
  \item A complete, documented, reproducible two-phase implementation of
        an autonomous greenhouse harvester, from world modelling to
        mission supervision, with every algorithm re-implemented rather
        than imported from a navigation stack.
  \item A framework-independent algorithm library validated in one
        simulator (Webots) and reused unchanged in another
        (Gazebo/ROS~2), demonstrating a practical portability pattern.
  \item A three-axis measurement layer --- pose accuracy, map accuracy,
        and time/throughput --- that scores the running system against
        simulator ground truth continuously, and which caught two
        regressions that no test and no visual inspection had caught.
  \item A pre-flight regression suite of 191 pure-Python checks, each
        one encoding a failure the project actually suffered, that runs
        in about one second before every simulation.
  \item An explicit SWOT analysis of every algorithm and architectural
        decision (Sec.~\ref{sec:swot}), including the ones that were
        abandoned.
\end{enumerate}

\subsection{How to read the flowcharts}

Every control-flow structure in this report is given as a flowchart
using one fixed visual language, defined once in
Fig.~\ref{fig:legend} and never redefined afterwards. Two conventions
matter:

\textbf{Shapes} follow classic flowchart practice, with two additions
specific to a robotics stack: a distinct shape for a \emph{message
read or written on a ROS topic}, and a distinct shape for a \emph{call
into the framework-independent algorithm library}. Being able to see, at
a glance, where a diagram touches the middleware and where it touches
the algorithms is the whole point of the split described in
Sec.~\ref{sec:architecture}.

\textbf{Phase bands} partition every diagram into three coloured
regions: \textcolor{phaseInitLine}{\textbf{INITIALISATION}} (executed
once, before any data arrives), \textcolor{phaseMainLine}{\textbf{MAIN
LOOP}} (the periodic callback, executed at a fixed rate for the whole
run), and \textcolor{phaseEndLine}{\textbf{TERMINATION / DESTRUCTION}}
(shutdown, resource release, final reporting). The banding is not
decoration: a recurring source of defects in this project was code
placed in the wrong band --- state armed once that should have been
re-armed per cycle, and cleanup that never ran because it was inside the
loop it was meant to end.

% =====================================================================
%  Fig. 1 -- the flowchart alphabet, defined once for the whole report.
%  Spans both columns (figure*) because it is a reference the reader
%  comes back to.
% =====================================================================
\begin{figure*}[!t]
\centering
\begin{tikzpicture}[fcchain]

% ---------------------------------------------------------------------
%  LEFT BLOCK : the six node shapes
% ---------------------------------------------------------------------
\node[fcterm] (s1) {START / STOP};
\node[fcproc, below=of s1] (s2) {Process\\{\tiny a computation}};
\node[fcdec,  below=of s2] (s3) {Decision?};
\node[fcio,   below=8mm of s3] (s4) {ROS topic in / out};
\node[fcsub,  below=of s4] (s5) {lib/ call\\{\tiny framework-independent}};
\node[fcfail, below=of s5] (s6) {Recovery / degraded\\{\tiny branch taken on failure}};

\node[font=\sffamily\footnotesize\bfseries, above=3mm of s1] (shapeshdr)
     {(a) Shapes};

% shape annotations, to the right of each
\node[anchor=west, font=\sffamily\scriptsize, right=3mm of s1]
     {Entry and exit points of the node.};
\node[anchor=west, font=\sffamily\scriptsize, right=3mm of s2]
     {Ordinary computation, no I/O.};
\node[anchor=west, font=\sffamily\scriptsize, right=3mm of s3]
     {Test; both exits are labelled \textsf{yes}/\textsf{no}.};
\node[anchor=west, font=\sffamily\scriptsize, right=3mm of s4]
     {A message published or received.};
\node[anchor=west, font=\sffamily\scriptsize, right=3mm of s5]
     {Call into \filename{lib/}: no ROS, no simulator.};
\node[anchor=west, font=\sffamily\scriptsize, right=3mm of s6]
     {Fallback path: brake, escape, abandon goal.};

% ---------------------------------------------------------------------
%  RIGHT BLOCK : the three phase bands
% ---------------------------------------------------------------------
\node[fcterm, right=78mm of s1] (p0) {node start};
\node[fcproc, below=of p0] (p1) {declare parameters,\\open publishers \& subscribers};
\node[fcproc, below=of p1] (p2) {arm timers, seed state};

\node[fcio,   below=9mm of p2] (p3) {wait for next message / tick};
\node[fcproc, below=of p3] (p4) {read inputs, run the algorithm};
\node[fcdec,  below=of p4] (p5) {shutdown\\requested?};
\node[fcio,   below=9mm of p5] (p6) {publish command};

\node[fcproc, below=10mm of p6] (p7) {stop the base, release the arm};
\node[fcproc, below=of p7] (p8) {emit final report};
\node[fcterm, below=of p8] (p9) {destroy node, rclpy.shutdown()};

\draw[fcline] (p0) -- (p1);
\draw[fcline] (p1) -- (p2);
\draw[fcline] (p2) -- (p3);
\draw[fcline] (p3) -- (p4);
\draw[fcline] (p4) -- (p5);
\draw[fcline] (p5) -- node[fclab,right] {no} (p6);
\draw[fcback] (p6.east) -- ++(14mm,0) |- (p3.east);
\node[font=\sffamily\tiny\itshape, text=phaseMainLine]
     at ($(p6.east)+(20mm,7mm)$) {loop};
\draw[fcline] (p5.west) -- ++(-13mm,0) |- node[fclab,pos=0.25] {yes} (p7.west);

\node[font=\sffamily\footnotesize\bfseries, above=3mm of p0] {(b) Phase bands};

% ---- the bands themselves, on the background layer -------------------
\begin{scope}[on background layer]
  \phaseband{bandinit}{INITIALISATION}{phaseInitLine}{(p0)(p1)(p2)}{bi}
  \phaseband{bandmain}{MAIN LOOP}{phaseMainLine}{(p3)(p4)(p5)(p6)}{bm}
  \phaseband{bandend}{TERMINATION / DESTRUCTION}{phaseEndLine}{(p7)(p8)(p9)}{be}
\end{scope}

% ---- band annotations ------------------------------------------------
\node[fcnote, right=6mm of bi.east, anchor=west, text width=44mm]
     {\textbf{Executed once.} Everything whose value must survive the
      whole run. Parameters are declared here so a YAML file can
      override them --- the source of one of the regressions reported in
      Sec.~\ref{sec:regression}.};
\node[fcnote, right=6mm of bm.east, anchor=west, text width=44mm]
     {\textbf{Executed at a fixed rate.} One tick, one decision. State
      that must be re-armed per cycle belongs here and not above: a
      counter armed once in initialisation is exactly the defect that
      disabled multi-pick harvesting (Sec.~\ref{sec:phaseB_mission}).};
\node[fcnote, right=6mm of be.east, anchor=west, text width=44mm]
     {\textbf{Executed once, on the way out.} Must run even on
      \texttt{SIGINT}: the performance monitor prints its final report
      from this band, because every run of this project ends with
      Ctrl\,+\,C.};

\end{tikzpicture}
\caption{The flowchart language used throughout this report, defined
once. (a) The six node shapes: two are specific to a robotics stack ---
a \emph{ROS topic} shape, so the boundary with the middleware is
visible, and a \emph{library call} shape, so the boundary with the
framework-independent algorithm layer is visible. (b) The three phase
bands into which every later flowchart is partitioned. The banding is
functional rather than decorative: two of the defects analysed in this
report are cases of code placed in the wrong band.}
\label{fig:legend}
\end{figure*}


\subsection{Report structure}

Sec.~\ref{sec:related} reviews the literature each algorithm is drawn
from. Sec.~\ref{sec:environment} specifies the toolchain, the exact
software versions, and the installation procedure.
Sections~\ref{sec:phaseA}--\ref{sec:phaseA_results} present \phaseA{}
(Webots) in full, subsystem by subsystem.
Sections~\ref{sec:phaseB}--\ref{sec:regression} present \phaseB{}
(ROS~2 / Gazebo) in the same order, so the two phases can be read side
by side. Sec.~\ref{sec:swot} gives the SWOT analysis of every method.
Sec.~\ref{sec:results} reports the measured results,
Sec.~\ref{sec:discussion} discusses the limitations and what they imply
for a physical deployment, and Sec.~\ref{sec:conclusion} concludes.
Appendix~\ref{app:reproduce} gives the commands to reproduce every
experiment; Appendix~\ref{app:parameters} tabulates every tuned
parameter with its justification; Appendix~\ref{app:topics} gives the
full topic contract; Appendix~\ref{app:troubleshooting} collects the
failure modes encountered and their diagnosis.
